\section{格林函数}
线性方程 $Lu=f$ 的格林函数是同一个算子在点源下的解,
\begin{equation}
  L G(\br,\brp)=\delta(\br-\brp),
\end{equation}
再配上和原问题相同的齐次边界条件.由于 $L$ 线性,源 $f$ 可以看成点源的叠加,
\begin{equation}
  u(\br)=\int G(\br,\brp)\,f(\brp)\,\mathrm{d}V'.
  \label{eq:green-superposition}
\end{equation}
前面热核\eqref{eq:heat-kernel}就是无限长杆上热方程的格林函数,三维波动方程的推迟解也是同一想法.

\subsection{一维两点边值}
在 $0<x<l$ 上求解 $-u''=f$,$u(0)=u(l)=0$.格林函数满足 $-G''=\delta(x-x')$,两端为零.点源两侧方程是齐次的,所以
\begin{equation}
  G=
  \begin{cases}
    A x, & 0\leq x\leq x', \\
    B(l-x), & x'\leq x\leq l.
  \end{cases}
\end{equation}
$G$ 在 $x'$ 连续,$A x'=B(l-x')$.把 $-G''=\delta$ 从 $x'-\epsilon$ 积到 $x'+\epsilon$,导数有一个跳跃,
\begin{equation}
  G'(x'+)-G'(x'-)=-1.
\end{equation}
左导数是 $A$,右导数是 $-B$,所以 $A+B=1$.于是
\begin{equation}
  G(x,x')=
  \begin{cases}
    x(l-x')/l, & x\leq x', \\
    x'(l-x)/l, & x\geq x'.
  \end{cases}
  \label{eq:green-1d}
\end{equation}
它对称,$G(x,x')=G(x',x)$.解是 $u(x)=\int_0^l G(x,x')f(x')\,\mathrm{d}x'$.

\begin{example}
$f=1$,即 $-u''=1$,$u(0)=u(l)=0$.
\end{example}
\begin{solution}
\begin{equation}
  u(x)=\int_0^x\frac{x'(l-x)}{l}\,\mathrm{d}x'
  +\int_x^l\frac{x(l-x')}{l}\,\mathrm{d}x'
  =\frac{x(l-x)}{2}.
\end{equation}
直接积分 $-u''=1$ 得到同一个抛物线,边界条件定掉两个常数.
\end{solution}

诺伊曼问题 $-u''=f$、$u'(0)=u'(l)=0$ 不能对任意 $f$ 求解.把方程从 $0$ 积到 $l$,左边是零,所以必须 $\int_0^l f=0$.满足这个条件时,解还差一个任意常数.第二类格林函数同样只能确定到相差一个常数,并且 $\partial G/\partial n$ 在边界上不能处处为零,而要等于 $-1/l$,用来和源的积分约束相配.

\subsection{格林公式}
设 $\varphi,\psi$ 在区域 $\Omega$ 上有足够的导数.散度定理用于 $\varphi\nabla\psi$ 和 $\psi\nabla\varphi$,相减得到格林第二恒等式
\begin{equation}
  \int_\Omega\bigl(\varphi\Delta\psi-\psi\Delta\varphi\bigr)\,\mathrm{d}V
  =\oint_{\partial\Omega}
    \Bigl(\varphi\funcpd{\psi}{n}-\psi\funcpd{\varphi}{n}\Bigr)\,\mathrm{d}S.
  \label{eq:green-identity}
\end{equation}
$\partial/\partial n$ 是外法向导数.

\subsection{泊松方程}
三维里 $-\Delta(1/R)=4\pi\delta^3(\br)$,$R=|\br|$.因此无界空间中 $-\Delta u=\rho$,并且 $u$ 在无穷远处趋于零时,
\begin{equation}
  u(\br)=\int\frac{\rho(\brp)}{4\pi|\br-\brp|}\,\mathrm{d}V'.
  \label{eq:poisson-free}
\end{equation}
$1/(4\pi R)$ 是 $-\Delta$ 的基本解.二维基本解是 $(1/(2\pi))\ln(1/R)$,因为 $\Delta\ln R=2\pi\delta^2(\br)$.

有界区域上的狄利克雷问题 $-\Delta u=\rho$,$u|_{\partial\Omega}=g$.取格林函数 $-\Delta G=\delta(\br-\brp)$,并且 $G$ 在 $\partial\Omega$ 上为零.把\eqref{eq:green-identity}用于 $u(\brp)$ 和 $G(\br,\brp)$,点源挑出 $u(\br)$,边界上 $G=0$,剩下
\begin{equation}
  u(\br)
  =\int_\Omega G(\br,\brp)\,\rho(\brp)\,\mathrm{d}V'
  -\oint_{\partial\Omega} g(\brp)\funcpd{G}{n'}\,\mathrm{d}S'.
  \label{eq:green-dirichlet}
\end{equation}
体积分是源的贡献,面积分是边界值的贡献.无界空间、边界在无穷远且 $g=0$ 时,它回到\eqref{eq:poisson-free}.

\subsection{电像法}
接地导体把边界电位定为零.格林函数等于基本解加上一个在区域内部调和的修正,这个修正在边界上正好抵消基本解.修正项如果能看成区域外面一个像电荷的电位,构造就变成电像法.

半空间 $z>0$,导体平面 $z=0$ 接地.单位点源放在 $(0,0,d)$,$d>0$.像电荷 $-1$ 放在 $(0,0,-d)$.于是
\begin{equation}
  G(\br,\brp)
  =\frac{1}{4\pi|\br-\brp|}
  -\frac{1}{4\pi|\br-\brp_*|},
  \qquad
  \brp_*=(x',y',-z').
\end{equation}
$z=0$ 时两个距离相等,$G=0$.像电荷在 $z>0$ 里是调和的,所以 $-\Delta G$ 在上半空间只剩下原来的点源.

接地导体球的半径是 $R$,球心在原点,单位点源在 $\brp$,$d=|\brp|>R$.像电荷放在同一射线上的
\begin{equation}
  \brp_*=\frac{R^2}{d^2}\brp,
\end{equation}
电荷量是 $-R/d$.球面上 $|\br|=R$ 时,余弦定理给出
\begin{equation}
  |\br-\brp_*|=\frac{R}{d}\,|\br-\brp|,
\end{equation}
所以
\begin{equation}
  \frac{1}{|\br-\brp|}-\frac{R/d}{|\br-\brp_*|}=0.
\end{equation}
格林函数是
\begin{equation}
  G(\br,\brp)
  =\frac{1}{4\pi|\br-\brp|}
  -\frac{R/d}{4\pi|\br-\brp_*|}.
\end{equation}
球外区域里只有第一个点源.球心处没有像电荷,$\brp_*$ 落在球内.

\begin{example}
单位点电荷放在接地球外,$d>R$.球面上的感应电荷总量.
\end{example}
\begin{solution}
无穷远处电位像 $1/(4\pi r)$ 减去 $(R/d)/(4\pi r)$.多出来的通量对应像电荷.感应电荷总量等于像电荷 $-R/d$.$d\to R^+$ 时它趋于 $-1$,点电荷几乎贴在球面上,感应电荷把它完全屏蔽.$d\to\infty$ 时像电荷趋于 $0$.
\end{solution}

\subsection{用本征函数展开}
有界区间上,格林函数也可以写成本征函数的级数.两端为零的热方程,$0<x<l$,
\begin{equation}
  G(x,t;x',t')
  =\frac{2}{l}\sum_{n=1}^{\infty}
    \sin\frac{n\pi x}{l}\sin\frac{n\pi x'}{l}
    \exp\Bigl(-a^2\bigl(\tfrac{n\pi}{l}\bigr)^2(t-t')\Bigr),
\end{equation}
只在 $t>t'$ 时使用,更早的时刻因果性要求 $G=0$.对 $x'$ 乘上当时的温度再积分,就是分离变量得到的解.级数在 $t=t'$ 时回到 $\delta(x-x')$,因为正弦族完备.
